\scanpage{190}

\[
\boxed{\boxed{\beta(A_1)=T}}\quad\text{i.e.}\quad A_1(\deleted{\unsure{1}}T)=T
\]

\deleted{\unsure{Posons|Donnons}} Considérons

\[
\deleted{\text{\unsure{Don|D'où|Donc}}}\qquad A_1-X_0=\varphi_0^0(X_0)-X_0=-\bigl(\varphi_0^0(X_1)-X_1\bigr)=f\in\mathcal A_1^0
\]

(où $\varphi_0(x)$ est la composante de $\varphi(x)$ suivant

\[
S\{X_0\}\simeq k\widetilde{\{T\}}\{X_0,X_1\}/(X_0+X_1-T)_{pd}\qquad\text{dans}
\]

\[
\mathcal A_1=\underbrace{S\{X_0\}}_{A_1^0}+\underbrace{S\{X_0\}dX_0}_{A_1^1}\simeq S\{X_1\}+S\{X_1\}dX_1\ \bigr)
\]

[ \deleted{NB} Comme $X_1=T-X_0$, on a

\[
\varphi_0^0(X_1)=\varphi_0^0(T)-\varphi_0^0(X_0)\inserted{=T-\varphi_0^0(X_0)},\quad\text{donc}\quad \varphi_0^0(X_1)-X_1=X_0-\varphi_0^0(X_0)\ \bigr]
\]

Alors la conjonction des conditions $A_1(X_0)=0$ et $A_1(T)=T$ s'écrit aussi

\[
 f\in(X_0)_{pd}\cap(X_1)_{pd}
\]

donc, si \unsure{l'on|on} interprète $\mathcal A_1^0$ comme un \deleted{\unsure{quotient de|produit de}} $k\{X_0,X_1\}$, cela signifie qu'on doit avoir

\[
 f=\sum_{i,j\geq\underline{\underline{1}}}\deleted{\unsure{a}}f_{ij}X_0^{(i)}X_1^{(j)}
 \qquad f_{ij}\in k
\]

d'où en résumé \deleted{$A_1=$}

\uline{Prop}. Les endom. $\varphi$ de la $S$-algèbre diff\textsuperscript{lle} $\mathcal A_1\simeq S\{X_0,X_1,dX_0,dX_1\}/(X_0+X_1-T,dX_0-dX_1)_{pd}$, (où $S=k\{T\}$ avec ses puiss. div) qui sont compatibles aux deux augm. $\alpha,\beta$ $\mathcal A_1\unsure{-|\to} S$ ($\alpha(X_0)=0,\alpha(X_1)=T$; $\beta(X_0)=T,\beta(X_1)=0$) sont donnés exactement par les systèmes

\[
 f\in\mathcal A_1^0=k\{X_0,X_1\}
\]

avec $f\in\operatorname{Ker}(k\{X_0,X_1\}\unsure{-|\to}k\{X_0\})\cap\operatorname{Ker}(k\{X_0,X_1\}\unsure{-|\leftarrow|\to}k\{X_1\})$

i.e. $\boxed{f=\sum_{i,j\geq1}a_{ij}X_0^{(i)}X_1^{(j)}}$

% archive sheet number: 189
